The fact that a compact operator is ``more than'' continuous motivated the use of the term \emph{completely continuous} operator for a compact operator. This terminology was common years ago but is used less frequently today.

The connection between compact operators and the dimension of the domain and range of the operator is even closer than Lemma~\ref{lem:infinite-dimensional-identity-noncompact} suggests.

\begin{Theorem}
\label{thm:finite-dimensional-domain-range-compactness}
Let $(\mathcal{D}(A),A)$ be a linear operator from $X$ to $Y$. Then we have the following:\begin{enumerate}\item If $(\mathcal{D}(A),A)$ is bounded and the range $\mathcal{R}(A)$ is finite-dimensional, then the operator $(\mathcal{D}(A),A)$ is compact.\item If the domain $\mathcal{D}(A)$ is finite-dimensional, then the operator $(\mathcal{D}(A),A)$ is compact.\end{enumerate}
\end{Theorem}

\begingroup
\postdisplaypenalty=10000
\renewcommand{\endProof}{\par\nobreak\hfill$\blacksquare$\par\vspace{0.1cm}}
\begin{Proof}
For part 1, let $x_n\in\mathcal{D}(A)$ be a given bounded sequence. Since the operator $(\mathcal{D}(A),A)$ is bounded, the sequence $A(x_n)\in\mathcal{R}(A)$ is also bounded. Since $\mathcal{R}(A)$ is finite-dimensional, the Bolzano-Weierstrass theorem implies that $A(x_n)$ has a convergent subsequence. Thus, $(\mathcal{D}(A),A)$ is compact.

For part 2, we note that the dimension of the range of an operator is less than or equal to the dimension of the domain. (If $\{x_i\}_{i=1}^k$ is a basis for $\mathcal{D}(A)$, then $\{A(x_i)\}_{i=1}^k$ spans $\mathcal{R}(A)$.) Also, by Lemma~\ref{lem:finite-dimensional-domain-operators-bounded}, any operator with a finite-dimensional domain is bounded. Thus, we can use part 1 to complete the proof.
\end{Proof}
\endgroup

\begin{Definition}
\label{def:finite-rank-operator}
If $A\in\mathcal{L}(X,Y)$ and $\mathcal{R}(A)$ is finite-dimensional, we say the operator $A$ has \textbf{finite rank}.
\end{Definition}

One common way of proving an operator is compact is by approximating by other operators (such as operators of finite rank) which are known to be compact. In using such an approximation scheme one usually employs the following result.

\begin{Theorem}
\label{thm:norm-limit-of-compact-operators}
Let $A_n\in\mathcal{L}(X,Y)$ be a sequence of compact operators. Suppose $A_n$ converges in the operator norm to an operator $A$. Then $A$ is compact.
\end{Theorem}

\begingroup
\postdisplaypenalty=10000
\renewcommand{\endProof}{\par\nobreak\hfill$\blacksquare$\par\vspace{0.1cm}}
\begin{Proof}
We employ a ``diagonal sequence'' argument. Let $\{x_n\}_{n=1}^\infty\subset X$ be a given bounded sequence. Then since $A_1$ is compact, the sequence $A_1(x_n)$ has a convergent subsequence. We label this subsequence $\{A_1(x_{1,n})\}_{n=1}^\infty$. Now, since $\{x_{1,n}\}_{n=1}^\infty$ is bounded and $A_2$ is compact, we see that $A_2(x_{1,n})$ has a convergent subsequence. We label this subsequence $\{A_2(x_{2,n})\}_{n=1}^\infty$. We now repeat the process, taking further subsequences of subsequences so that $\{x_{k,n}\}_{n=1}^\infty$ is a subsequence of $\{x_{j,n}\}_{n=1}^\infty$ if $j<k$ and so that $\{A_k(x_{k,n})\}_{n=1}^\infty$ converges. (Recall that since $\{A_k(x_{k,n})\}_{n=1}^\infty$ is convergent it is Cauchy.)

Now consider the diagonal sequence $\{x_{n,n}\}_{n=1}^\infty$. We denote $z_n:=x_{n,n}$. Note that this is indeed a subsequence of the original sequence $x_n$. We claim that $A(z_n)$ is Cauchy and hence convergent. (This will complete the proof since $x_n$ was an arbitrary bounded sequence.) Let $\varepsilon>0$ be given. We note that for any $i$, $j$ and $k$, we have
\begin{equation}\label{eq:compact-operator-limit-three-term-estimate}
\begin{aligned}
\|A(z_i)-A(z_j)\|&\leq\|A(z_i)-A_k(z_i)\|+\|A_k(z_i)-A_k(z_j)\|\\
&\quad+\|A(z_j)-A_k(z_j)\|.
\end{aligned}
\end{equation}
Since $z_n$ is a bounded sequence, and since $A_k\to A$ in the operator norm, we can pick $k$ sufficiently large so that
\begin{equation}\label{eq:compact-operator-limit-uniform-tail-estimate}
\|A(z_n)-A_k(z_n)\|<\varepsilon/3
\end{equation}
for every element of the sequence $z_n$. We now note that for fixed $k$, the sequence $A_k(z_n)$ is Cauchy. This is true since $\{z_n\}_{n=k}^\infty$ is a subsequence of $\{x_{k,m}\}_{m=1}^\infty$. Thus, we can pick $i$ and $j$ sufficiently large so that
\begin{equation}\label{eq:compact-operator-limit-subsequence-cauchy-estimate}
\|A_k(z_i)-A_k(z_j)\|<\varepsilon/3.
\end{equation}
Combining \eqref{eq:compact-operator-limit-three-term-estimate} with \eqref{eq:compact-operator-limit-uniform-tail-estimate} and \eqref{eq:compact-operator-limit-subsequence-cauchy-estimate} completes the proof.
\end{Proof}
\endgroup
